% ================================================================
% ==== FILE: content/cycles/cycles_k3.tex
% ================================================================

\subsubsection{Zyklen auf \texorpdfstring{$K_3$}{K_3}}
\label{sec:cycles-k3}

Ebene$K_3$entspricht dem RAF-chemischen Kontinuum:
ein sich selbst erhaltendes Reaktionsnetzwerk, das durch Autokatalyse ermöglicht wird und
F-Verschluss.
Sein Zustandsraum ist gegeben durch
\[
  \Omega(K_3) = \{(M,R,C,F)\mid (M',R')\in \mathrm{RAF}(G_2)\},
\]
Wo$G_2$ist das Konnektivitätsdiagramm, von dem geerbt wurde$K_2$,
und die RAF-Bedingung erzwingt sowohl autokatalytische Unterstützung als auch Erreichbarkeit
aller Reaktanten aus dem Lebensmittelset$F$.
Zyklen weiter$K_3$entstehen als wiederkehrende chemische Transformationen
Reaktionswege, katalytische Flüsse und Wiederherstellung des Verschlusses.

\subsubsection{Grundlegende Arten von Zyklen auf \texorpdfstring{$K_3$}{K_3}}

Drei strukturelle Klassen von Zyklen definieren die dynamische Komplexität von
das RAF-Kontinuum.

\paragraph{(1) Autocatalytic Reaction Cycles.}

Dies sind geschlossene Pfade im Reaktionsraum:
\[
  m_{i_1}\xrightarrow{r_1} m_{i_2}
    \xrightarrow{r_2} \dots
    \xrightarrow{r_k} m_{i_1},
\]
where each reaction $r_j$ is catalysed by some $m\in M'$.
Solche Zyklen halten Konzentrationen aufrecht, kontrollieren Produktionsverhältnisse und
propagieren RAF-Stabilität.
Sie stellen das minimale dynamische Substrat dar$K_3$.

\paragraph{(2) F-closure Restoration Cycles.}

Die RAF-Eigenschaft erfordert:
\[
  \text{all reactants in } R' \text{ must be reachable from } F.
\]
Perturbations (e.g.\ deficit of a molecule) can temporarily break
F-Verschluss.
Zyklen, die den F-Abschluss wiederherstellen, erfüllen:
\[
  F \to M_{\mathrm{lost}} \to M_{\mathrm{restored}} \to F
\]
und sind für die Lebensfähigkeit von wesentlicher Bedeutung$K_3$.
They interact directly with the threshold $\Theta_{\mathrm{kas}}$.

\paragraph{(3) Energetic and Redox Cycles.}

Following the extended chemical block (the corresponding derivation note),
Zu den chemischen Potenzialen gehören:
\[
  P_{\mathrm{chem}}(t),\;
  P_{\mathrm{redox}}(t),\;
  P_{\mathrm{grad}}(t).
\]
Entsprechend,$K_3$unterstützt energetische Kreisläufe:
\[
  C_{\mathrm{energy}}:
  P_{\mathrm{chem}}\to P_{\mathrm{redox}}
    \to P_{\mathrm{grad}} \to P_{\mathrm{chem}},
\]
die die Stoffwechsel- und Gradientenzyklen von vorwegnehmen$K_4$.

Diese Prozesse halten das RAF-System im energieeffizienten Bereich
des Zustandsraums.

\subsubsection{Strukturelle Spannungen und Schwellen}

Die strukturelle Spannung$T_3(t)$wird bestimmt durch:
\[
T_3 =
   f\bigl(
     \mathrm{deficit}_{\mathrm{catalysis}},\;
     1-\rho_{\mathrm{closure}},\;
     \mathrm{energy\ deficit},\;
     \mathrm{redox\ imbalance}
   \bigr),
\]
where $\rho_{\mathrm{closure}}$ is the closure metric defined by:
\[
  \rho_{\mathrm{closure}} =
    \frac{\text{number of reachable reactants from } F}
         {\text{total reactants in } R'}.
\]

Die Lebensfähigkeit des Zyklus erfordert:
\[
  T_3(t) < \Theta_{\mathrm{death}}^{(3)},
\]
und die Existenz einer RAF erfordert:
\[
  \rho_{\mathrm{cat}} \cdot \rho_{\mathrm{closure}}
  \ge \Theta_{\mathrm{kas}}.
\]

So geht es weiter$K_3$existieren innerhalb einer klar definierten Lebensfähigkeitsregion.

\subsubsection{Metriken von Zyklen auf \texorpdfstring{$K_3$}{K_3}}

Im Einklang mit den universellen Metriken (the corresponding construction route):

\paragraph{Length.}
\[
  L(C) = \int_0^\tau \|J_{\mathrm{chem}}(t)\|\, dt,
\]
where $J_{\mathrm{chem}}$ is the reaction–flow vector.

\paragraph{Efficiency.}
\[
  C_{\mathrm{eff}}
   = \frac{1}{L(C)}
     \int_0^\tau J_{\mathrm{chem}}(t)\cdot dA_{\mathrm{chem}}(t),
\]
where $A_{\mathrm{chem}}$ are chemical-state axes (active/inactive
Reaktion, An-/Abwesenheit von Arten).

\paragraph{Stability.}
\[
  S(C) = \min_{t\in[0,\tau]}
         d\bigl(
           \Omega(K_3),\;
           \partial\Omega(K_3)
         \bigr),
\]
was der Nähe zum Brechen entweder der Autokatalyse oder entspricht
Schließung.

\paragraph{Weight.}
\[
  w(C) = F\bigl(
           C_{\mathrm{eff}},\;
           S(C),\;
           \mathrm{energy\ turnover}
         \bigr).
\]

\subsubsection{Zeit- und Zyklusperioden}

Die chemische Zeit läuft$K_3$ist definiert über:
\[
  \tau(K_3) = \int \frac{ds}{v_{\mathrm{chem}}},
\]
where $v_{\mathrm{chem}}$ is the effective reaction velocity.
Für RAF-Zyklen beträgt die charakteristische Wiederholungsperiode:
\[
  \tau_{\mathrm{RAF}}
  \sim \frac{1}{k_{\mathrm{cat}}}
\]
where $k_{\mathrm{cat}}$ is the effective catalytic rate.

Die Zeit selbst ist somit eine Funktion der wiederkehrenden Reaktionsstruktur.

\subsubsection{Geburt höherer Zyklen}

Zyklen bei$K_3$dienen als Vorläufer$K_4$Zyklen:

\begin{itemize}
  \item Membranwartungszyklen,
  \item Energie-Gradienten-Zyklen über eine Grenze hinweg,
  \item protometabolische Schleifen,
  \item Vorerregbarkeitszyklen (über Ionen- oder Protonengradienten).
\end{itemize}

The continuity map $K_3\to K_4$ (the corresponding derivation note)
zeigt, dass Zyklen von:

\[
  (P_{\mathrm{chem}}, P_{\mathrm{grad}}, J_{\mathrm{chem}})
\]

nach und nach werden:

\[
  (P_{\mathrm{mem}}, P_{\mathrm{grad}}, J_{\mathrm{in/out}})
\]

as a membrane $\partial\Omega$ forms.

So geht es weiter$K_3$Kodieren Sie die früheste Organisationsstruktur.
Biologische Zyklen (Stoffwechsel, Membran, Erregbarkeit) entstehen erst später
von dieser Basis aus.
